\documentclass[12pt,a4paper]{article}

\usepackage[brazil]{babel}
\usepackage{fontspec}
\usepackage{xeCJK}
\usepackage[table,xcdraw]{xcolor}
\usepackage{amsmath}
\usepackage{float}
\usepackage{tabularx}
\usepackage{array}
\usepackage{hyperref}
\usepackage{tabularx}
\usepackage{seqsplit}


\setmainfont{Latin Modern Roman}
\setCJKmainfont{Noto Serif CJK JP}

\begin{document}

\section*{Projeto: 5ª Oficina de Origami Seishun}

\subsection*{Informações Gerais}

\begin{table}[H]
\centering
\begin{tabularx}{\textwidth}{|>{\raggedright\arraybackslash}p{5cm}|X|}
\hline
\textbf{Informação} & \textbf{Descrição} \\ \hline

Período de Realização &
01 de Outubro \\ \hline

Carga Horária Total &
2 horas \\ \hline

Horário &
19h - 21h \\ \hline

Local &
UTFPR - Sala:A 022 \\ \hline

Quantidade de Vagas &
25 \\ \hline

Custo &
R\$ 0,00 \\ \hline

Público-alvo &
Pessoas com idade entre 5 e 50 anos \\ \hline

Perfil dos Participantes &
Pessoas interessadas em aprender origami e cultura japonesa \\ \hline

Pré-requisitos &
Nenhum \\ \hline

Divulgação &
Comunidade interna e externa, por meio das redes sociais \\ \hline

\end{tabularx}
\end{table}

\subsection*{Objetivo Geral}

Proporcionar aos participantes uma introdução à arte do origami, desenvolvendo
habilidades de coordenação motora, concentração e criatividade, enquanto
aprendem a criar diversas figuras em papel.

\subsection*{Conteúdo Programático}

Nossa oficina está estruturada em três ``Ilhas'' principais, cada uma
correspondendo a um nível de dificuldade progressivo. Isso permite que todos,
desde iniciantes até participantes com alguma experiência, possam aproveitar
a oficina de acordo com seu nível de habilidade.

\subsection*{Metodologia}

\paragraph{Ilhas do Origami:}
Haverá três ``ilhas'' disponíveis. Cada ilha representará um nível de
dificuldade diferente: fácil, médio e difícil. Isso permite que cada
participante escolha a ilha que melhor se adapta à sua habilidade ou à sua
vontade de aprender.

\paragraph{Monitores Dedicados:}
Em cada ilha, haverá de 1 a 3 monitores prontos para oferecer suporte
individualizado. Eles estarão disponíveis para tirar dúvidas, demonstrar
técnicas e auxiliar em qualquer dificuldade que surgir.

\paragraph{Passo a Passo Visual:}
Em cada ilha, os participantes terão acesso a instruções visuais para a
realização de cada origami. Esses guias passo a passo facilitarão o
entendimento do processo de dobradura.

\paragraph{Material para Treino:}
Para que os participantes possam praticar à vontade e adquirir confiança,
será disponibilizado papel sulfite em todas as ilhas. A proposta é permitir
que cada participante experimente quantas vezes forem necessárias.

\paragraph{Seu Origami para Casa:}
Ao final da atividade, cada participante receberá um papel próprio para
realizar sua peça final e poderá levá-la para casa.

\subsection*{Critérios de Avaliação}

\paragraph{Desenvolvimento de Habilidades de Dobradura:}
Será observado se o participante conseguiu realizar as dobras com destreza
e seguir corretamente os passos descritos.

\paragraph{Resolução de Problemas e Persistência:}
Será observado como o participante reage aos obstáculos encontrados durante
a atividade e sua capacidade de superá-los.

\subsection*{Instrutores}



\begin{table}[H]
\centering
\scriptsize
\setlength{\tabcolsep}{3pt}

\begin{tabularx}{\textwidth}{|>{\raggedright\arraybackslash}X|>{\centering\arraybackslash}p{2.3cm}|>{\raggedright\arraybackslash}X|>{\centering\arraybackslash}p{2.4cm}|}
\hline
\rowcolor[HTML]{E0E0E0} 
\textbf{Nome} & \textbf{CPF} & \textbf{E-mail} & \textbf{Telefone} \\ \hline

André Luiz da Silva Junior & 841.696.098-4 & \seqsplit{and695842@gmail.com} & (41) 98830-8688 \\ \hline
Ayumi Wassano & 549.396.978-51 & \seqsplit{ayumiwassano@alunos.utfpr.edu.br} & (14) 99879-3615 \\ \hline
Bianca Natsumi Nishimura & 430.971.558-30 & \seqsplit{biancanishimura@alunos.utfpr.edu.br} & (11) 99561-9524 \\ \hline
Brian Yuta Maruyama Tatewaki & 063.945.809-29 & \seqsplit{brian.tatewaki@select.eng.br} & (43) 99185-1775 \\ \hline
Gustavo Arai Matsubara & 112.052.939-57 & \seqsplit{gustavoarai@alunos.utfpr.edu.br} & (43) 99130-1149 \\ \hline
João Pedro da Silva & 141.801.669-17 & \seqsplit{jsilva.2007@alunos.utfpr.edu.br} & (43) 9955-3680 \\ \hline
João Vitor Vieira Ribeiro Doneda & 507.694.748-88 & \seqsplit{joaodoneda@alunos.utfpr.edu.br} & (11) 99811-7281 \\ \hline
Thalita Yukare Pereira Kinoshita & 476.568.448-21 & \seqsplit{thalitayukare@alunos.utfpr.edu.br} & (11) 96572-5541 \\ \hline
\end{tabularx}
\end{table}

\subsection*{Equipe de Apoio}

\begin{table}[H]
\centering
\scriptsize
\setlength{\tabcolsep}{3pt}
\begin{tabularx}{\textwidth}{|>{\raggedright\arraybackslash}X|>{\centering\arraybackslash}p{2.3cm}|>{\raggedright\arraybackslash}X|>{\centering\arraybackslash}p{2.4cm}|}
\hline
\rowcolor[HTML]{E0E0E0} 
\textbf{Nome} & \textbf{CPF} & \textbf{E-mail} & \textbf{Telefone} \\ \hline
Anna Liz Michlovská R. Rodrigues & 098.944.436.80 & \seqsplit{annarodrigues.2006@alunos.utfpr.edu.br} & (64) 99926-7884 \\ \hline
Camilla Kaori Oshiro & 119.473.519-31 & \seqsplit{camillaoshiro@alunos.utfpr.edu.br} & (43) 99115-1267 \\ \hline
Douglas Spigolon Barbosa & 409.729.398-22 & \seqsplit{douglas\_sbarbosa@hotmail.com} & (14) 99737-3902 \\ \hline
Ghabriel Jun Aizawa & 666.560.994-4 & \seqsplit{ghabriellaizawa@gmail.com} & (43) 98461-0584 \\ \hline
Gustavo Bashio Nakano & 387.428.198-14 & \seqsplit{gustavo.bashio.nakano@gmail.com} & (17) 99112-9740 \\ \hline
João Pedro Diniz Nacur & 100.997.349-58 & \seqsplit{luisoujpof@gmail.com} & (43) 99930-0814 \\ \hline
Júlio Martins Bertonha & 448.002.208-23 & \seqsplit{jbertonha@alunos.utfpr.edu.br} & (19) 99773-0331 \\ \hline
Luana Ayumi Nishiyama da Silva & 098.901.979-95 & \seqsplit{luanas.1999@utfpr.alunos.edu.br} & (43) 99962-4350 \\ \hline
Lucas Santos de Oliveira & 105.011.139-75 & \seqsplit{lucasoliveira.2023@alunos.utfpr.edu.br} & (43) 99672-7165 \\ \hline
Melissa Mayumi Hukuchima & 424.260.368-18 & \seqsplit{melissahukuchima@alunos.utfpr.edu.br} & (14) 99172-1138 \\ \hline
Rebeca Sayuri Ogawa & 121.499.089-48 & \seqsplit{rebecasayuriogawa@alunos.utfpr.edu.br} & (43) 98829-4700 \\ \hline
Renato André Antunes Kanieski & 094.602.949-01 & \seqsplit{renatokanieski@alunos.utfpr.edu.br} & (42) 99134-8894 \\ \hline

\end{tabularx}
\end{table}

\end{document}